\documentclass[11pt]{article}
\usepackage{fontspec}
\setmainfont{Latin Modern Roman}
\usepackage{amsmath,amssymb,mathtools,unicode-math}
\setmathfont{Latin Modern Math}
\usepackage{xurl,hyperref}
\hypersetup{colorlinks=true,urlcolor=blue,pdftitle={Problem statement}}
\newcommand{\sref}[2]{\hyperlink{source:#1:#2}{[S#2]}}
\newcommand{\cataloguescope}[1]{\paragraph{Mathematical question.}#1}
\begin{document}
\section{Research Problems in Function Theory — Problem 3.16}\label{statement}

\subsection{Definitions and mathematical statement}

\paragraph{Shared notation.}$\mathbb N=\{1,2,\ldots\}$, and $\mathbb Z,\mathbb Q,\mathbb R,\mathbb C$ denote the integers, rationals, reals and complex numbers. Any other conventions are specified locally.


For compact $K\subset\mathbb R^n$, $n\geq3$, define Newtonian capacity by
\[\operatorname{cap}(K)=\left(\inf_{\mu\in\mathcal P(K)}\iint |x-y|^{2-n}\,d\mu(x)\,d\mu(y)\right)^{-1},\]
where $\mathcal P(K)$ denotes probability measures supported on $K$, $1/\infty=0$, and $\operatorname{cap}(\varnothing)=0$. A set is polar if each of its compact subsets has capacity zero. For a compact $E$ set $c(P,r)=\operatorname{cap}(E\cap\overline B(P,r))$ and
\[T(E)=\left\{P\in E:\int_0^1\frac{c(P,r)}{r^{n-1}}\,dr<\infty\right\}.\]
\paragraph{Statement recorded in the supplied source.}
Can one prove directly, without invoking Kellogg's theorem, that $T(E)$ is polar for every compact $E\subset\mathbb R^n$, and thereby clarify the sharpness of the displayed integrated Wiener criterion? \sref{statement}{2}
\subsection{Short English statement}
Give a direct proof that the points where a compact set satisfies the thinness integral form a polar set, preferably explaining why the integral condition is sharp.
\subsection{Sources}
\begin{description}
\item[\hypertarget{source:statement:1}{S1}] ULAM.AI and the UnsolvedMath Contributors, UnsolvedMath: A Curated Collection of Open Mathematics Problems, catalogue entry. CC BY 4.0. \url{https://huggingface.co/datasets/ulamai/UnsolvedMath}.
\item[\hypertarget{source:statement:2}{S2}] W. K. Hayman and E. F. Lingham, \emph{Research Problems in Function Theory}, arXiv:1809.07200v2 (2018), Problem 3.16. \url{https://arxiv.org/abs/1809.07200}.
\end{description}
\label{end:statement}
\end{document}
